%%%%%%%%%%%%%%%%%%%%%%%%%%%%%%%%%%%%%%%%%%%%%%%%%%%%%%%%%%%%%%%%%%%%%%%%%%%%%%%%
% Medium Length Graduate Curriculum Vitae
% LaTeX Template
% Version 1.2 (3/28/15)
%
% This template has been downloaded from:
% http://www.LaTeXTemplates.com
%
% Original author:
% Rensselaer Polytechnic Institute 
% (http://www.rpi.edu/dept/arc/training/latex/resumes/)
%
% Modified by:
% Daniel L Marks <xleafr@gmail.com> 3/28/2015
%
% Important note:
% This template requires the res.cls file to be in the same directory as the
% .tex file. The res.cls file provides the resume style used for structuring the
% document.
%
%%%%%%%%%%%%%%%%%%%%%%%%%%%%%%%%%%%%%%%%%%%%%%%%%%%%%%%%%%%%%%%%%%%%%%%%%%%%%%%%

%-------------------------------------------------------------------------------
%	PACKAGES AND OTHER DOCUMENT CONFIGURATIONS
%-------------------------------------------------------------------------------

%%%%%%%%%%%%%%%%%%%%%%%%%%%%%%%%%%%%%%%%%%%%%%%%%%%%%%%%%%%%%%%%%%%%%%%%%%%%%%%%
% You can have multiple style options the legal options ones are:
%
%   centered:	the name and address are centered at the top of the page 
%				(default)
%
%   line:		the name is the left with a horizontal line then the address to
%				the right
%
%   overlapped:	the section titles overlap the body text (default)
%
%   margin:		the section titles are to the left of the body text
%	~/Documents/trislaz.github.io/_draf
%   11pt:		use 11 point fonts instead of 10 point fonts
%
%   12pt:		use 12 point fonts instead of 10 point fonts
%
%%%%%%%%%%%%%%%%%%%%%%%%%%%%%%%%%%%%%%%%%%%%%%%%%%%%%%%%%%%%%%%%%%%%%%%%%%%%%%%%

\documentclass[margin]{res}  

% Default font is the helvetica postscript font
\usepackage{helvet}
\usepackage{hyperref}

% Increase text height
\textheight=700pt

\begin{document}

%-------------------------------------------------------------------------------
%	NAME AND ADDRESS SECTION
%-------------------------------------------------------------------------------
\par
\name{Tristan Lazard}

% Note that addresses can be used for other contact information:
% -phone numbers
% -email addresses
% -linked-in profile

\address{37 Boulevard Mortier\\75020, Paris\\\href{https://twitter.com/trislaz}{@trislaz}}
\address{+336.89.95.84.82\\trislaz.tl@gmail.com\\ \href{trislaz.github.io}{trislaz.github.io}}

% Uncomment to add a third address
%\address{Address 3 line 1\\Address 3 line 2\\Address 3 line 3}
%-------------------------------------------------------------------------------

\begin{resume}

%-------------------------------------------------------------------------------
%	EDUCATION SECTION
%-------------------------------------------------------------------------------


\section{EDUCATION}
\textbf{PSL University}, Paris, 75006\\
{\sl ITI - MS Innovation }, Entrepreneurship and innovation. \hfill 2018-2019.
\par
\textbf{Sorbonne University}, Paris  \hfill 2016-2017\\
{\sl Master 2 : Applied mathematics}
Thesis: \href{https://drive.google.com/file/d/16VLpiMhV8CecNDiYf85nuo1kFdp9lQAF/view?usp=sharing}{"A model of random mutualistic network."}. 
\par
\textbf{Ecole Normale Supérieure}, rue d'Ulm, Paris\\
{\sl Master 1 : Evolutionary biology and ecology} \hfill 2015-2016 \\
{\sl Bachelor in life-sciences} \hfill 2013-2014
\par 
\textbf{Lycée Saint-Louis}, Paris \hfill 2011-2013\\
{\sl Classe préparatoire BCPST} 
\par 
\textbf{AgroCampus-Ouest}, Rennes \hfill 2014-2015\\
{\sl Bachelor in Agronomy, spare year} 
%-------------------------------------------------------------------------------

%-------------------------------------------------------------------------------


%-------------------------------------------------------------------------------
%	COMPUTER SKILLS SECTION
%-------------------------------------------------------------------------------

\section{SKILLS}
\textbf{Languages}: French (mother tongue), fluent English and debutant in Spanish.
\\
\textbf{Programming languages}: Python, Bash, \LaTeX, nextflow. 
\\
\textbf{Computing skills}: docker, slurm, cloud-computing, GPU-computing, 
ML+image processing usual libs (pytorch, openCV, PIL, sklearn...).
\\

%-------------------------------------------------------------------------------

%-------------------------------------------------------------------------------
%	EXPERIENCE SECTION
%-------------------------------------------------------------------------------
% Modify the format of each position
\begin{format}
\title{l}\employer{r}\\
\dates{l}\location{r}\\
\body\\
\end{format}
%-------------------------------------------------------------------------------
\par
\section{PROFESSIONAL\\EXPERIENCE}
\employer{CBIO - Mines-Paristech}
\location{Paris, France}
\dates{December 2019 - December 2023}
\title{\textbf{PhD in BioInformatics}}
\begin{position}
My PhD project focuses on processing Whole Slide Images, which are digitized microscopic images of diseased tissue commonly used by clinicians for diagnosis. These images contain a vast amount of information, leading to large file sizes (10-15 GB uncompressed). By developing machine learning algorithms to predict patient-level labels from these images, we can assist clinicians in their practice and improve our understanding of the disease through the interpretation of the algorithm's decisions. Main Achievements:
\begin{itemize}
    \item \textbf{5 publications}, 3 first authors, in biological journals and machine learning venues. 
    \item 1 patent.
    \item \textbf{Prize-winning solution (5000\$)} at the \href{https://www.drivendata.org/competitions/148/visiomel-melanoma/leaderboard/}{VisioMel challenge}. 
        \href{https://github.com/drivendataorg/visiomel-melanoma/tree/main/3rd%20Place}{Link to the solution}
\end{itemize}

\end{position}

\employer{CMM, Mines-Paristech}
\location{Fontainebleau, France}
\dates{Jan - July 2019}
\title{\textbf{Engineering internship}}
\begin{position}
{\sl Automatic counting of Melanocytes}: In partnership with L'Oréal,
I created an algorithm for automatically counting melanocytes in microscopic images. This solution has been integrated into the processing pipeline and is now utilized in regular research.
\end{position}
\\
\employer{IBPS, Sorbonne university}
\location{Paris, France}
\dates{Jan - August 2018}
\title{\textbf{Research internship}}
\begin{position}
{\sl Large scale phylogenetic tree inference}: Develop an algorithm for calculating phylogenetic trees from huge genetic sequence data (-omics data).
.
\end{position}

%-------------------------------------------------------------------------------

\section{PUBLICATIONS}
\textbf{Lazard, T.}, Bataillon, G., Naylor, P., Popova, T.,
Bidard, G.-C., Stoppa-Lyonnet, D., Stern, M.-
H., Decencie` re, E., Walter, T., and Vincent-Sal-
omon, A. (2022). Deep Learning identifies new
morphological patterns of Homologous
Recombination Deficiency in luminal breast
cancers from whole slide images. \textbf{Cell Reports
Medicine} 
\par
\textbf{Lazard, T.}, Lerousseau, M., Decenci\`ere, E., Walter, T. (2023).
Giga-SSL: Self-Supervised Learning for Gigapixel Images.
Proceedings of the IEEE/CVF Conference on Computer Vision and Pattern Recognition \textbf{(CVPR) Workshops}.
\par
\textbf{Lazard, T.}, Blusseau S., Velasco-Forero S, Decenci\`ere E., Flouret V., Cohen C., Baldeweck T.
(2022). Applying Deep Learning to Melanocyte Counting on Fluorescent TRP1 Labelled Images of In Vitro Skin Model.
\textbf{Image Analysis & Stereology}.
\par
Lubrano M., \textbf{Lazard T.}, Balezo G., Bellahsen-Harrar Y., Badoual C., Berlemont S., Walter T. (2022).
\textbf{Computer Vision–ECCV} 2022 Workshops: Tel Aviv, Israel, October 23–27, 2022, Proceedings.
\par
Naylor P, \textbf{Lazard T}, Bataillon G, Laé M, Vincent-Salomon A, Hamy AS, et al. Prediction of Treatment
Response in Triple Negative Breast Cancer From Whole Slide Images. \textbf{Front Signal Process}.
\par
\textbf{Lazard T.}, Decenci\`ere E., Walter T. Democratizing computational pathology: optimized WSI representations for the TCGA.
\textbf{Preprint}.
\par
\textbf{T. Lazard}, et. al., Cancer du sein - Utilisation de l’intelligence artificielle pour prédire le statut tumoral relatif à la recombinaison homologue (2023). \textbf{Médecine et Sciences}
\par
A. Beaufr\`ere*, \textbf{T. Lazard*}, et. al. Self-supervised learning to predict intrahepatic cholangiocarcinoma transcriptomic classes on routine histology (2024). \textbf{Preprint}
\\

%-------------------------------------------------------------------------------
%	Interests
%-------------------------------------------------------------------------------
\section{INTERESTS}
Climbing, harmonica, poetry and litterature. 
%-------------------------------------------------------------------------------
\end{resume}
\end{document}
